\subsection{Progent：可编程动作权限与 Utility-Security Tradeoff}

slides 76、78--85 聚焦 Progent。Motivating examples 展示宽泛工具权限如何让恶意或误导指令触发高风险动作；Progent 把自然语言/程序化 policy 转成 runtime privilege control，在 Agent 生成 action 时检查用户目标、资源、参数和上下文。评测关注 attack success rate 的下降，同时检查正常任务 utility 是否保留。

\lecturefigure{slides-images/slide-076.png}{官方 slide 76：Progent motivating example 起点}
\lecturefigure{slides-images/slide-078.png}{官方 slide 78：高权限 Agent 的完整风险示例}
\lecturefigure{slides-images/slide-079.png}{官方 slide 79：Progent privilege-control overview}
\lecturefigure{slides-images/slide-082.png}{官方 slide 82：Policy 与 action validation 的完整流程}
\lecturefigure{slides-images/slide-083.png}{官方 slide 83：Progent runtime enforcement 组件}
\lecturefigure{slides-images/slide-084.png}{官方 slide 84：Progent 显著降低 attack success}
\lecturefigure{slides-images/slide-085.png}{官方 slide 85：安全增益与任务 utility 的联合评测}

\begin{importantbox}{Policy Enforcement 的三项输入}
一次 action 至少要结合：用户已经授权的目标、当前 Agent/工具拥有的 capability、以及资源与参数的具体上下文。仅按 tool name allowlist 不够，因为同一邮件工具可以发送正常通知，也可以外泄数据；同一文件工具可以读取项目目录，也可以读取 secret。
\end{importantbox}

\begin{knowledgebox}{读图：安全评测必须同时画两条曲线}
Progent 的价值不是把 ASR 降到零而让所有任务失败，而是在保持 utility 的同时压缩攻击面。比较方案时应报告正常任务成功率、延迟、人工审批负担和攻击成功率；若只报安全指标，最安全的“关闭 Agent”会错误地成为最优解。
\end{knowledgebox}

